\section{模空间的离壳描述}
\label{sec:off-shell-description-moduli-space}
\setcounter{equation}{2}

我们希望把 $J$-全纯映射
$u:(\Sigma,j)\to(M,J)$ 看成某个巴拿赫向量丛的弗雷德霍姆截面在
$\mathcal F:=C^\infty(\Sigma,M)$ 上的一个零点。

我们先假设 $\Sigma$ 没有边界，稍后再指出处理有边情形所需的改动。

$C^\infty(\Sigma,M)$ 关于 $C^\infty$ 拓扑实际上并不是巴拿赫流形，而只是弗雷歇流形。为了从 $C^\infty(\Sigma,M)$ 得到巴拿赫流形，需要取 $C^\infty(\Sigma,M)$ 关于 $W^{k,p}$ 范数（相应地，关于 $C^{k,\alpha}$ 赫尔德范数）的完备化，将它扩大为 $k-2/p>0$ 时的索伯列夫映射空间
$\mathcal F^{k,p}:=W^{k,p}(\Sigma,M)$（相应地，赫尔德空间
$C^{k,\alpha}(\Sigma,M)$）。我们将只使用索伯列夫空间，并先回顾取值于光滑流形 $M$ 的 $W^{k,p}$ 映射概念。写出其定义最容易也最快捷的办法，是使用纳什等距嵌入定理 (Nas56)。

固定一个等距嵌入
$(M,g)\hookrightarrow\mathbb R^N$，其中整数 $N$ 充分大，并把 $M$ 看成
$\mathbb R^N$ 的闭子流形。于是可以把
$u:\Sigma\to M\hookrightarrow\mathbb R^N$ 看成 $\mathbb R^N$ 值函数。

\begin{definition}\label{def:sobolev-maps-manifold}
设 $k-2/p\geq0$。定义
\[
 W^{k,p}(\Sigma,M):=
 \{u\in W^{k,p}(\Sigma,\mathbb R^N)\mid u(z)\in M\ \text{a.e.}\},
\]
并把 $W^{k,p}(\Sigma,M)$ 的任意元素 $u$ 称为到 $M$ 的 $W^{k,p}$ 映射。
\end{definition}

容易验证，只要取 $k-2/p>0$，从而映射 $u$ 连续，这个定义就与嵌入无关，因而在流形上是良定义的。

\begin{remark}\label{rem:local-definition-sobolev-maps}
也可以局部定义有关范数，再在 $M$ 的一个固定局部有限图册上利用单位分解求和，由此定义完备化。在实际应用中，尤其是研究映射在 $W^{k,p}$ 拓扑中的收敛时，需要使用先验估计方法，此时这个局部定义会更实用。
\end{remark}

现在设
$u\in C^\infty(\Sigma,M)\subset W^{k,p}(\Sigma,M)$。考虑集合
$W^{k,p}(u^*TM)$，其定义为
\[
 W^{k,p}(u^*TM)=C^\infty(u^*TM)\ \text{关于 $W^{k,p}$ 的完备化}.
\]
回忆 $C^\infty(u^*TM)$ 定义为
\[
 C^\infty(u^*TM)=
 \{\xi:\Sigma\to TM\mid\pi\circ\xi=u,\ \xi\ \text{光滑}\},
\]
其中 $\pi:TM\to M$ 是显然的投影。我们也用
$\pi:u^*TM\to\Sigma$ 表示诱导投影
\[
 u^*TM|_z=T_{u(z)}M\longrightarrow z,
 \qquad z\in\Sigma.
\]

以 $\mathcal F^{k,p}=W^{k,p}(\Sigma,M)$ 表示 $\mathcal F$ 的 $W^{k,p}$ 完备化。

\begin{proposition}\label{prop:sobolev-mapping-space-banach-manifold}
假设 $M$ 是有限维闭流形，并设 $k-2/p>0$。则
$\mathcal F^{k,p}$ 带有 $C^\infty$ 巴拿赫流形结构，并且它在 $u$ 处的切空间为
\begin{equation}
 T_u\mathcal F^{k,p}=W^{k,p}(u^*TM).
 \label{eq:tangent-sobolev-mapping-space}
\end{equation}
\end{proposition}

\begin{proof}
需要为 $\mathcal F^{k,p}$ 给出一个图册。固定 $M$ 上的度量 $g$，并令
$\delta=\iota(M,g)$ 为 $(M,g)$ 的单射半径。记
\[
 D^\delta(TM)=\{(x,v)\in TM\mid |v|<\delta\}.
\]
于是相伴指数映射
\[
 \operatorname{Exp}:D_\delta(M)\subset TM\longrightarrow M\times M
\]
是良定义的，并且是到其像上的微分同胚。记
$\operatorname{Exp}(x,v)=(x,\exp_x(v))$。考虑以\emph{光滑映射 $u$ 为中心}的开子集
\[
 \mathcal U_u=\widetilde{\exp}_u\mathcal V_u,
\]
其中
\[
 \mathcal V_u=
 \{\xi\in W^{k,p}(u^*TM)\mid |\xi(z)|<\iota(M,g)\},
\]
而 $\widetilde{\exp}_u:\mathcal V_u\to\mathcal U_u$ 由公式
\[
 \widetilde{\exp}_u(\xi)(z):=\exp_{u(z)}\xi(z)
\]
定义。于是 $\widetilde{\exp}_u$ 是到其像上的同胚，而该像包含 $u$；因此
$(\mathcal U_u,\widetilde{\exp}_u^{-1})$ 定义了
$W^{k,p}(\Sigma,M)$ 的一个局部坐标卡，其模型巴拿赫空间为
$W^{k,p}(u^*TM)$。

在 $W^{k,p}(\Sigma,M)$ 中选取一个由光滑映射组成的可数稠密子集，并把这一构造用于其中每个元素。这样的选取来自嵌入
$C^\infty(\Sigma,M)\hookrightarrow W^{k,p}(\Sigma,M)$ 的稠密性，于是得到
\[
 W^{k,p}(\Sigma,M)=
 \bigcup_{u\in C^\infty(\Sigma,M)}\mathcal U_u.
\]
接下来证明这些局部坐标卡的相容性。
\end{proof}

\begin{lemma}\label{lem:exponential-charts-smooth-atlas}
集合族
\[
 \{\widetilde{\exp}_u^{-1}:\mathcal U_u\to W^{k,p}(u^*TM)
 \mid u\in C^\infty(\Sigma,M)\}
\]
构成 $\mathcal F^{k,p}$ 的一个 $C^\infty$ 图册。
\end{lemma}

\begin{proof}
需要证明：对任意选择的
$u_1,u_2\in C^\infty(\Sigma,M)$，转移映射
\[
 \widetilde{\exp}_{u_1}^{-1}\circ\widetilde{\exp}_{u_2}:
 \widetilde{\exp}_{u_2}^{-1}(\mathcal U_{u_2}\cap\mathcal U_{u_1})
 \longrightarrow
 \widetilde{\exp}_{u_1}^{-1}(\mathcal U_{u_2}\cap\mathcal U_{u_1})
\]
对所有 $\ell\in\mathbb N$ 都是 $C^\ell$ 的。

\begin{figure}[htbp]
 \centering
 \begin{tikzpicture}[x=0.78cm,y=0.78cm,>=Stealth,line cap=round,line join=round]
  \begin{scope}[shift={(0,3.45)}]
    \begin{scope}
      \clip
        (-2.25,0.85) .. controls (-2.95,1.15) and (-3.15,0.45) .. (-2.63,0.18)
        .. controls (-2.05,-0.13) and (-2.48,-0.80) .. (-2.12,-1.12)
        .. controls (-1.53,-1.62) and (-0.65,-1.16) .. (-0.45,-0.60)
        .. controls (-0.20,-0.10) and (-0.30,0.60) .. (-0.75,0.93)
        .. controls (-1.23,1.16) and (-1.87,1.01) .. cycle;
      \clip
        (-0.24,1.23) .. controls (0.55,1.73) and (1.55,1.32) .. (1.92,0.57)
        .. controls (2.28,-0.18) and (2.04,-1.07) .. (1.34,-1.36)
        .. controls (0.70,-1.63) and (-0.10,-1.43) .. (-0.50,-0.91)
        .. controls (-0.85,-0.45) and (-0.92,0.76) .. (-0.24,1.23);
      \fill[gray!45] (-3.2,-1.8) rectangle (2.5,1.8);
    \end{scope}
    \draw
      (-2.25,0.85) .. controls (-2.95,1.15) and (-3.15,0.45) .. (-2.63,0.18)
      .. controls (-2.05,-0.13) and (-2.48,-0.80) .. (-2.12,-1.12)
      .. controls (-1.53,-1.62) and (-0.65,-1.16) .. (-0.45,-0.60)
      .. controls (-0.20,-0.10) and (-0.30,0.60) .. (-0.75,0.93)
      .. controls (-1.23,1.16) and (-1.87,1.01) .. cycle;
    \draw
      (-0.24,1.23) .. controls (0.55,1.73) and (1.55,1.32) .. (1.92,0.57)
      .. controls (2.28,-0.18) and (2.04,-1.07) .. (1.34,-1.36)
      .. controls (0.70,-1.63) and (-0.10,-1.43) .. (-0.50,-0.91)
      .. controls (-0.85,-0.45) and (-0.92,0.76) .. (-0.24,1.23);
  \end{scope}

  \begin{scope}[shift={(-3.0,0)}]
    \clip (0,0) circle (1.45);
    \fill[gray!45] (0.92,0) ellipse [x radius=0.58,y radius=1.02];
  \end{scope}
  \draw (-3,0) circle (1.45);
  \draw (-4.85,0)--(-1.15,0) (-3,-1.85)--(-3,1.85);

  \begin{scope}[shift={(3.0,0)}]
    \clip (0,0) circle (1.45);
    \fill[gray!45] (-0.92,0) ellipse [x radius=0.58,y radius=1.02];
  \end{scope}
  \draw (3,0) circle (1.45);
  \draw (1.15,0)--(4.85,0) (3,-1.85)--(3,1.85);

  \draw[->] (-3.65,1.80) .. controls (-3.60,2.45) and (-2.93,2.83) .. (-2.20,2.89);
  \node[above] at (-3.05,2.40) {$\widetilde{\exp}_{u_1}$};
  \draw[->] (3.65,1.80) .. controls (3.60,2.45) and (2.93,2.83) .. (2.20,2.89);
  \node[above] at (3.05,2.40) {$\widetilde{\exp}_{u_2}$};
\end{tikzpicture}

 \caption{指数坐标卡。}
 \label{fig:exponential-chart}
\end{figure}

由 $\widetilde{\exp}$ 的定义，
\[
 (\widetilde{\exp}_{u_1}^{-1}\circ
 \widetilde{\exp}_{u_2})(\xi)(z)
 =\exp_{u_1(z)}^{-1}\circ\exp_{u_2(z)}(\xi(z)).
\]
注意，由于所选 $u_1,u_2$ 光滑，映射
$z\mapsto\widetilde{\exp}_{u_1}^{-1}\circ
\widetilde{\exp}_{u_2}(\xi(z))$ 属于 $W^{k,p}$ 当且仅当
$\xi$ 属于 $W^{k,p}$，因而它定义了 $u_1^*TM$ 的一个 $W^{k,p}$ 截面。此外，再次由 $u_1,u_2$ 的光滑性，赋值
$\xi\mapsto\widetilde{\exp}_{u_1}^{-1}\circ
\widetilde{\exp}_{u_2}$ 定义了一个映射
\begin{equation}
 \mathcal V_{u_1}\cap
 \widetilde{\exp}_{u_1}^{-1}\circ\widetilde{\exp}_{u_2}(\mathcal V_2)
 \longrightarrow
 \mathcal V_{u_2}\cap
 \widetilde{\exp}_{u_2}^{-1}\circ\widetilde{\exp}_{u_1}(\mathcal V_1),
 \label{eq:exponential-chart-transition-map}
\end{equation}
它无穷多次可微。交换 $u_1,u_2$ 的角色，也证明了逆映射
$\widetilde{\exp}_{u_2}^{-1}\circ\widetilde{\exp}_{u_1}$ 无穷多次可微。因此已经证明转移映射是 $C^\infty$ 的。见图~\ref{fig:exponential-chart}。
\end{proof}

\begin{exercise}\label{ex:smoothness-exponential-transition-map}
证明转移映射\eqref{eq:exponential-chart-transition-map}无穷多次可微，并证明关于切空间的陈述。
\end{exercise}

这就完成了证明。

\begin{remark}\label{rem:varying-model-banach-spaces}
模型巴拿赫空间 $W^{k,p}(u^*TM)$ 随基点 $u$ 的选择而变；但只要
$u_1,u_2$ 在 $C^\infty$ 意义下彼此充分接近，它们就是同胚的。因此，可以再复合一个从
$W^{k,p}(u_2^*TM)$ 到 $W^{k,p}(u_1^*TM)$ 的同胚，把命题~\ref{prop:sobolev-mapping-space-banach-manifold}所给的图册改造成 $\mathcal F^{k,p}$ 的一个图册，使其中的转移映射无穷多次可微，从而得到定义~\ref{def:banach-manifold}意义下的巴拿赫流形结构。
\end{remark}

从现在起，除非为了清楚起见绝对有必要，我们都在记号中略去上标 $k,p$。

把赋值 $u\mapsto du$ 看成某个（无穷维）向量丛的截面，下面说明这一点。记
\[
 \mathcal H_u=\Gamma(\Lambda^1(u^*TM)),
 \qquad
 \mathcal H=\bigcup_{u\in\mathcal F}\mathcal H_u.
\]
于是映射 $d:u\mapsto du$ 定义了丛 $\mathcal H\to\mathcal F$ 的一个截面。不过，这个截面不是弗雷德霍姆截面：注意，向量丛 $u^*TM$ 的秩为 $2n$，而
$T\Sigma\otimes u^*TM$ 的秩为 $4n$。依照标准记号，记
\[
 \Gamma(u^*TM)=\Omega^0(u^*TM),
 \qquad
 \Gamma(\Lambda^1(u^*TM))=\Omega^1(u^*TM).
\]

接下来引入 $\Sigma$ 上给定的复结构 $j$ 以及 $(M,\omega)$ 上的殆复结构 $J$。数对 $(j,J)$ 诱导分解
\[
 \operatorname{Hom}(T\Sigma,u^*TM)
 =\operatorname{Hom}'_{(j,J)}(T\Sigma,u^*TM)
 \oplus
 \operatorname{Hom}''_{(j,J)}(T\Sigma,u^*TM)
\]
为复线性部分和复反线性部分，即
\[
 A=\frac{A-J\cdot A\cdot j}{2}
   +\frac{A+J\cdot A\cdot j}{2}
   =:A'+A''.
\]
等价地，把 $A$ 看成 $\Lambda^1(u^*TM)$ 的元素，也可以关于殆复结构 $j,J$ 把它分解为
\[
 A=A^{(1,0)}+A^{(0,1)}.
\]
这诱导 $u^*TM$ 值 $1$-形式的相应分解
\[
 \Omega^1(u^*TM)
 =\Omega^{(1,0)}_{(j,J)}(u^*TM)
 \oplus\Omega^{(0,1)}_{(j,J)}(u^*TM),
\]
以及
\[
 \mathcal H=\mathcal H^{(0,1)}\oplus\mathcal H^{(0,1)},
\]
其中
\[
 \mathcal H^{(1,0)}_{((j,u),J)}
 =\Omega^{(1,0)}_{(j,J)}(u^*TM),
 \qquad
 \mathcal H^{(0,1)}_{((j,u),J)}
 =\Omega^{(0,1)}_{(j,J)}(u^*TM).
\]

最后定义\emph{截面}
\[
 \overline\partial_{(j,J)}:\mathcal F\longrightarrow
 \mathcal H^{(0,1)}_{(j,J)}
 :=\mathcal H^{0,1}_{j,J}(\mathord\cdot^*TM)
\]
为
\[
 u\longmapsto
 \overline\partial_{(j,J)}u:=\frac{du+J\cdot du\cdot j}{2}.
\]
由定义，从 $(\Sigma,j)$ 到 $(M,J)$ 的 $(j,J)$-全纯映射就是截面
$\overline\partial_{(j,J)}$ 的零点。以后在不致混淆时，我们将略去 $\Sigma$ 上的 $j$。还要指出
\[
 \operatorname{rank}u^*TM
 =\operatorname{rank}\Lambda^{(0,1)}_{(j,J)}(u^*TM)(=2n),
\]
这反映了 $\Sigma$ 的二维性。在不致混淆时，我们也将从记号中略去 $(j,J)$。
